\begin{corollary}[Analytic continuation and functional equation for the Riemann zeta function] \label{corollary:analytic_continuation_and_functional_equation_for_riemann_zeta_function}
\TODO{AI generated, verify}
The \CrefAndHyperrefIfExist{definition:riemann_zeta_function}{Riemann zeta function} $\zeta(s)$ admits a meromorphic continuation to $\bbC$ with a single simple pole at $s = 1$ of residue $1$.

Define the \hldef{completed Riemann zeta function}
$$
\boxed{ \xi(s) = \frac{1}{2} s (s - 1) \pi^{-s/2} \Gamma\!\left( \frac{s}{2} \right) \zeta(s). }
$$
Then $\xi(s)$ is an entire function satisfying the symmetric functional equation
$$
\boxed{ \xi(s) = \xi(1 - s). }
$$

Equivalently, in terms of the \CrefAndHyperrefIfExist{definition:riemann_zeta_function}{Riemann zeta function}:
$$
\zeta(s) = 2^s \pi^{s-1} \sin\!\left( \frac{\pi s}{2} \right) \Gamma(1-s) \, \zeta(1-s), \qquad s \in \bbC \setminus \{1\}.
$$

\begin{proof}
	\TODO{AI generated, verify}
	We specialize \Cref{theorem:analytic_continuation_and_functional_equation_for_dirichlet_L_function} to the trivial Dirichlet character modulo $1$. Let $\chi_0$ denote the principal character modulo $k = 1$, so that $\chi_0(n) = 1$ for all $n \geq 1$. The associated Dirichlet $L$-function is
	$$L(s, \chi_0) = \sum_{n=1}^\infty \frac{1}{n^s} = \zeta(s).$$
	\CrefAndHyperrefIfExist{definition:dirichlet_L_function_of_a_dirichlet_character_modulo_a_nonnegative_integer}{Dirichlet $L$-function}

	Since $\chi_0(-1) = 1$, the character is even, so $a(\chi_0) = 0$. The Gauss sum and root number are
	$$\tau(\chi_0) = \sum_{m=1}^1 \chi_0(m) e^{2\pi i m / 1} = 1,$$
	and
	$$\varepsilon(\chi_0) = \frac{\tau(\chi_0)}{i^{a(\chi_0)} \sqrt{1}} = 1.$$

	The completed Dirichlet $L$-function for $\chi_0$ coincides with the completed Riemann zeta function:
	$$\Lambda(s, \chi_0) = \left( \frac{1}{\pi} \right)^{s/2} \Gamma\!\left( \frac{s}{2} \right) \zeta(s) = \frac{1}{\frac{1}{2} s (s - 1)} \, \xi(s).$$
	The functional equation $\Lambda(s, \chi_0) = \varepsilon(\chi_0) \, \Lambda(1 - s, \overline{\chi_0})$ simplifies to
	$$\xi(s) = \xi(1 - s),$$
	because $\overline{\chi_0} = \chi_0$ and $\varepsilon(\chi_0) = 1$.

	We verify that $\xi(s)$ is entire. The only singularity of $\zeta(s)$ in $\bbC$ is a simple pole at $s = 1$ with residue $1$\CrefAndHyperrefIfExist{definition:riemann_zeta_function}{Riemann zeta function}. At $s = 1$, the factor $s - 1$ in $\xi(s)$ cancels this pole, so that
	$$\lim_{s \to 1} \xi(s) = \frac{1}{2} \cdot 1 \cdot \pi^{-1/2} \Gamma\!\left( \frac{1}{2} \right) \cdot 1 = \frac{1}{2},$$
	where we used $\Gamma(1/2) = \sqrt{\pi}$\Cref{proposition:value_of_gamma_function_at_one_half}. The functions $s(s-1)$, $\pi^{-s/2}$, and $\Gamma(s/2)$ are holomorphic for all $s \in \bbC$ except possibly at the poles of $\Gamma(s/2)$, which occur at non-positive even integers; at those points $\zeta(s)$ vanishes, so the product remains finite. Therefore, $\xi(s)$ extends to an entire function on $\bbC$.

	We derive the classical functional equation for $\zeta(s)$ from the identity $\xi(s) = \xi(1 - s)$. Substituting the definition of $\xi$ and dividing both sides by $s(s-1)$, we obtain
	$$\pi^{-s/2} \Gamma\!\left( \frac{s}{2} \right) \zeta(s) = \pi^{-(1-s)/2} \Gamma\!\left( \frac{1 - s}{2} \right) \zeta(1 - s).$$
	Solving for $\zeta(s)$ gives
	$$\zeta(s) = \pi^{s-1/2} \, \frac{\Gamma((1-s)/2)}{\Gamma(s/2)} \, \zeta(1 - s).$$
	To convert this to the classical form we use two standard identities for the Gamma function. Euler's reflection formula\Cref{proposition:eulers_reflection_formula_for_gamma_function} states that $\Gamma(z)\Gamma(1-z) = \pi / \sin(\pi z)$. Setting $z = (s+1)/2$ and using $\sin(\pi(s+1)/2) = \cos(\pi s/2)$, we obtain
	$$\Gamma\!\left( \frac{s+1}{2} \right) \Gamma\!\left( \frac{1-s}{2} \right) = \frac{\pi}{\cos(\pi s / 2)}.$$
	Legendre's duplication formula\Cref{proposition:legendres_duplication_formula_for_gamma_function} states that $\Gamma(z)\Gamma(z+1/2) = 2^{1-2z}\sqrt{\pi}\,\Gamma(2z)$. Setting $z = s/2$, we obtain
	$$\Gamma\!\left( \frac{s}{2} \right) \Gamma\!\left( \frac{s+1}{2} \right) = 2^{1-s} \sqrt{\pi} \, \Gamma(s).$$
	Combining these two identities, we can express $\Gamma((1-s)/2)$ in terms of $\Gamma(s/2)$ and $\Gamma(s)$: from the reflection formula, $\Gamma((s+1)/2) = \frac{\pi}{\cos(\pi s/2)\,\Gamma((1-s)/2)}$, and substituting into the duplication formula gives
	$$\Gamma\!\left( \frac{s}{2} \right) \cdot \frac{\pi}{\cos(\pi s / 2) \, \Gamma((1-s)/2)} = 2^{1-s} \sqrt{\pi} \, \Gamma(s).$$
	Rearranging yields
	$$\frac{\Gamma((1-s)/2)}{\Gamma(s/2)} = \frac{\pi}{2^{1-s}\sqrt{\pi}\,\Gamma(s)\cos(\pi s / 2)}.$$
	Applying Euler's reflection formula once more with $z = s$, we have $\Gamma(s) = \frac{\pi}{\sin(\pi s)\,\Gamma(1-s)}$. Substituting and using $\sin(\pi s) = 2\sin(\pi s/2)\cos(\pi s/2)$:
	$$\frac{\Gamma((1-s)/2)}{\Gamma(s/2)} = \frac{\pi}{2^{1-s}\sqrt{\pi} \cdot \frac{\pi}{\sin(\pi s)\,\Gamma(1-s)} \cdot \cos(\pi s / 2)} = \frac{2^s}{\sqrt{\pi}} \, \frac{\sin(\pi s/2)\cos(\pi s/2)}{\cos(\pi s/2)} \, \Gamma(1-s) = \frac{2^s}{\sqrt{\pi}} \sin\!\left( \frac{\pi s}{2} \right) \Gamma(1-s).$$
	Substituting this ratio back into the expression for $\zeta(s)$:
	$$\zeta(s) = \pi^{s-1/2} \cdot \frac{2^s}{\sqrt{\pi}} \sin\!\left( \frac{\pi s}{2} \right) \Gamma(1-s) \, \zeta(1-s) = 2^s \pi^{s-1} \sin\!\left( \frac{\pi s}{2} \right) \Gamma(1-s) \, \zeta(1-s).$$
	This is the classical functional equation.

	\CrefAndHyperrefIfExist{definition:riemann_zeta_function}{Riemann zeta function}
	\Cref{theorem:analytic_continuation_and_functional_equation_for_dirichlet_L_function}
\end{proof}
\end{corollary}